% Alıştırma bankası — bölüm 4
% Bölüm, bu dosyanın adına göre belirlenir.

\begin{exo}[Bir diferansiyel formu bir yol boyunca tümlemek]{integrale-forme-differentielle}
\keywords{diferansiyel form, tam diferansiyel, Schwarz ölçütü, çizgi integrali}

Şu iki diferansiyel formu ele alalım:
\[
\omega_1 = y\,\mathrm{d}x + x\,\mathrm{d}y,
\qquad
\omega_2 = y\,\mathrm{d}x,
\]
Ayrıca başlangıç noktası $O(0,0)$ ile $M(1,1)$ noktasını birleştiren üç yol
tanımlayalım: $\gamma_1$, yatay $O\to(1,0)$ parçasının ardından dikey
$(1,0)\to M$ parçasından; $\gamma_2$, dikey $O\to(0,1)$ parçasının ardından
yatay $(0,1)\to M$ parçasından oluşsun; $\gamma_3$ ise $O$ ile $M$'yi doğrudan
birleştiren doğru parçası olsun.

\begin{enumerate}
\item $\omega_1=\mathrm{d}f$ olacak biçimde bir $f$ fonksiyonu bularak $\omega_1$'in tam olduğunu gösterin.
\item Her parçayı parametreleyerek $\int_{\gamma_1}\omega_1$ ve $\int_{\gamma_2}\omega_1$ integrallerini hesaplayın. Sonucu hesap yapmadan yeniden elde edin.
\item Schwarz ölçütünü $\omega_2$'ye uygulayın. Ne sonuç çıkarılabilir?
\item $\int_{\gamma_1}\omega_2$, $\int_{\gamma_2}\omega_2$ ve $\int_{\gamma_3}\omega_2$ integrallerini hesaplayıp yorumlayın.
\end{enumerate}

\begin{indication}
Yatay bir parçada $\mathrm{d}y=0$, dikey bir parçada $\mathrm{d}x=0$'dır;
dolayısıyla her integral basit bir integrale indirgenir.
\end{indication}

\begin{solution}
\textbf{Soru 1.} $\partial f/\partial x=y$ ve $\partial f/\partial y=x$ olduğundan $f(x,y)=xy$ uygundur.

\textbf{Soru 2.} $t\mapsto\bigl(x(t),y(t)\bigr)$ ile parametrelenmiş bir eğri için
\[
\int_\gamma\omega_1
= \int \left[y(t)x'(t)+x(t)y'(t)\right]\,\mathrm{d}t.
\]

$\gamma_1$ yolu iki parçanın birleşimidir. İlk parçada $t\in[0,1]$ için
$x(t)=t$, $y(t)=0$; dolayısıyla $\mathrm{d}x=\mathrm{d}t$ ve
$\mathrm{d}y=0$'dır. İkinci parçada yine $t\in[0,1]$ için $x(t)=1$,
$y(t)=t$; dolayısıyla $\mathrm{d}x=0$ ve $\mathrm{d}y=\mathrm{d}t$'dır. Böylece
\[
\begin{aligned}
I_1
= \int_{\gamma_1}\omega_1
&= \int_0^1\left(0\times 1+t\times 0\right)\,\mathrm{d}t
 + \int_0^1\left(t\times 0+1\times 1\right)\,\mathrm{d}t \\
&= 0+\int_0^1\mathrm{d}t = 1.
\end{aligned}
\]

$\gamma_2$ için ilk parça $x(t)=0$, $y(t)=t$, ikinci parça ise
$x(t)=t$, $y(t)=1$ ile parametrelenir; her iki durumda da $t\in[0,1]$'dir.
Bu nedenle
\[
\begin{aligned}
I_2
= \int_{\gamma_2}\omega_1
&= \int_0^1\left(t\times 0+0\times 1\right)\,\mathrm{d}t
 + \int_0^1\left(1\times 1+t\times 0\right)\,\mathrm{d}t \\
&= 0+\int_0^1\mathrm{d}t = 1.
\end{aligned}
\]
İki değer aynıdır. Bu sonuç hesap yapmadan da öngörülebilirdi: $f(x,y)=xy$
olmak üzere $\omega_1=\mathrm{d}f$ olduğundan integral yalnız uç noktalara bağlıdır:
\[
\int_\gamma\omega_1=f(M)-f(O)=f(1,1)-f(0,0)=1.
\]

\textbf{Soru 3.} $\omega_2=A(x,y)\,\mathrm{d}x+B(x,y)\,\mathrm{d}y$ yazalım.
Burada $A(x,y)=y$ ve $B(x,y)=0$'dır. Form tam olsaydı Schwarz ölçütü
\[
\frac{\partial A}{\partial y}=\frac{\partial B}{\partial x}.
\]
eşitliğini verirdi. Oysa $\partial A/\partial y=1$, buna karşılık
$\partial B/\partial x=0$'dır. Çapraz türevler farklı olduğundan $\omega_2$
tam değildir. Bu nedenle integrali $O$ ile $M$ arasında izlenen yola bağlı
olabilir; aşağıdaki hesap bunu doğrular.

\textbf{Soru 4.} $\omega_2=y\,\mathrm{d}x$ olduğundan integrale yalnızca
$y$ sıfırdan farklıyken gerçekleşen bir $x$ değişimi katkıda bulunabilir.
Önceki parametrelemelerle $\gamma_1$ üzerinde
\[
\begin{aligned}
J_1
=\int_{\gamma_1}\omega_2
&=\int_0^1 0\times 1\,\mathrm{d}t
  +\int_0^1 t\times 0\,\mathrm{d}t \\
&=0.
\end{aligned}
\]
$\gamma_2$'nin ilk parçasında $x(t)=0$, dolayısıyla $\mathrm{d}x=0$;
ikinci parçasında $x(t)=t$, $y(t)=1$ ve $\mathrm{d}x=\mathrm{d}t$'dır. Böylece
\[
\begin{aligned}
J_2
=\int_{\gamma_2}\omega_2
&=\int_0^1 t\times 0\,\mathrm{d}t
  +\int_0^1 1\times 1\,\mathrm{d}t \\
&=1.
\end{aligned}
\]
Son olarak köşegen $\gamma_3$ açıkça
\[
x(t)=t,\qquad y(t)=t,\qquad t\in[0,1],
\]
ile parametrelenir; buradan $\mathrm{d}x=\mathrm{d}t$ ve
$\mathrm{d}y=\mathrm{d}t$ olur. Sonuç
\[
J_3
=\int_{\gamma_3}\omega_2
=\int_0^1 y(t)x'(t)\,\mathrm{d}t
=\int_0^1 t\,\mathrm{d}t
=\frac12.
\]
Üç yol aynı uç noktalara sahipken üç farklı değer verir: $J_1=0$, $J_2=1$
ve $J_3=1/2$. Bu, tam olmayan bir diferansiyel formun ayırt edici yol bağımlılığıdır.
\end{solution}
\end{exo}

\begin{exo}[Aynı iç enerji değişimi, üç aktarım biçimi]{meme-delta-u-trois-transferts}
\keywords{birinci yasa, iç enerji, ısı, mekanik iş, elektrik işi, kalorimetri, durum fonksiyonu}

Bir sıvı, kalorimetresi ve sıvıya daldırılmış küçük bir direnç, makroskopik
olarak durgun kapalı bir sistem oluşturmaktadır. Kap rijittir ve sistemin toplam
ısı kapasitesinin sabit olduğu varsayılmaktadır:
\[
C=2{,}00\ \mathrm{kJ\,K^{-1}}.
\]
Sistem, $T_A=20{,}0\,^\circ\mathrm{C}$ sıcaklığındaki $A$ denge durumundan
$T_B=25{,}0\,^\circ\mathrm{C}$ sıcaklığındaki $B$ denge durumuna getirilecektir.
Şu bağıntı kabul edilsin:
\[
U(B)-U(A)=C(T_B-T_A).
\]
Makroskopik kinetik ve potansiyel enerji değişimleri sıfırdır. Dışarıya olan
tüm kayıplar ihmal edilmektedir.

Aynı başlangıç durumu $A$'dan başlayarak aşağıdaki üç protokol ayrı ayrı uygulanır:
\begin{itemize}
\item[Protokol 1] Sistem, kendisine hiçbir iş verilmeden sıcak bir kaynakla ısıl temasa sokulur.
\item[Protokol 2] Duvarlar ısı yalıtımlıdır. Bir palet, $m$ kütlesinin $h=5{,}00$ m yükseklikten düşmesi sayesinde sıvıyı karıştırır. Palet ve sıvı sonunda durur; kütlenin yitirdiği tüm mekanik enerji sisteme aktarılır. $g=9{,}81\ \mathrm{m\,s^{-2}}$ alınacaktır.
\item[Protokol 3] Sistem $Q_3=4{,}00$ kJ ısı alır ve kalan enerji dirence elektrik yoluyla verilir. Alınan elektrik gücü sabit ve $\mathcal P=300$ W'tır.
\end{itemize}

\begin{enumerate}
\item Üç protokol için ortak olan $\Delta U=U(B)-U(A)$ iç enerji değişimini hesaplayın.
\item İlk iki protokolün her biri için sistemin aldığı ısı $Q$ ile iş $W$'yi belirleyin. İkinci protokolde gerekli $m$ kütlesini hesaplayın.
\item Üçüncü protokolde alınan elektrik işini ve direncin beslenme süresini hesaplayın.
\item Üç süreç aynı başlangıç ve son durumlarına sahiptir. Buna rağmen $Q$ ve $W$ değerlerinin neden farklı olabileceğini açıklayın. Sistem son $B$ durumunda “ısı” ya da “iş” içerir mi?
\end{enumerate}

\begin{indication}
Bankacı işaret uzlaşımıyla $\Delta U=Q+W$ bağıntısını uygulayın. Palet için
alınan iş, kütlenin potansiyel enerjisindeki $mgh$ azalmaya eşittir. Tellerden
sınırı geçen elektrik enerjisi elektrik işi olarak sayılır.
\end{indication}

\begin{solution}
\textbf{Soru 1.} Sıcaklık farkı $T_B-T_A=5{,}0$ K'dir. Dolayısıyla
\[
\Delta U=C(T_B-T_A)
=2{,}00\times5{,}0
=10{,}0\ \mathrm{kJ}.
\]
Bu değer yalnızca $A$ ve $B$ denge durumlarına bağlıdır.

\textbf{Soru 2.} Birinci protokolde iş alınmaz: $W_1=0$. Birinci yasa
\[
Q_1=\Delta U=10{,}0\ \mathrm{kJ}.
\]
sonucunu verir. Enerji sisteme girdiği için ısı pozitiftir.

İkinci protokolde duvarlar ısı yalıtımlı olduğundan $Q_2=0$'dır. İç enerji
değişiminin tamamı paletin mekanik işinden kaynaklanır:
\[
W_2=\Delta U=10{,}0\ \mathrm{kJ}.
\]
$W_2=mgh$ olduğuna göre
\[
m=\frac{W_2}{gh}
=\frac{10{,}0\times10^3}{9{,}81\times5{,}00}
\simeq 204\ \mathrm{kg}.
\]
Bu büyük mertebe, küçük bir sıcaklık artışının bile önemli bir mekanik
enerjiye karşılık geldiğini hatırlatır.

\textbf{Soru 3.} Birinci yasa gereğince
\[
W_3=\Delta U-Q_3=10{,}0-4{,}00=6{,}00\ \mathrm{kJ}.
\]
Bu iş pozitiftir: elektrik beslemesi sisteme enerji verir. $W_3=\mathcal P\,\Delta t$ ile
\[
\Delta t=\frac{W_3}{\mathcal P}
=\frac{6{,}00\times10^3}{300}
=20{,}0\ \mathrm{s}.
\]

\textbf{Soru 4.} Üç yol sırasıyla
\[
(Q_1,W_1)=(10{,}0,0)\ \mathrm{kJ},
\qquad
(Q_2,W_2)=(0,10{,}0)\ \mathrm{kJ},
\]
ve
\[
(Q_3,W_3)=(4{,}00,6{,}00)\ \mathrm{kJ}.
\]
sonuçlarını verir. Her durumda $Q+W$ toplamı $10{,}0$ kJ'dir ve aynı
$\Delta U$ değişimini üretir. İç enerji bir durum fonksiyonudur; ısı ve iş ise
süreç sırasındaki aktarımları betimler. Son $B$ durumunda sistem iç enerjiye
sahiptir, fakat “ısı” ya da “iş” içermez.
\end{solution}
\end{exo}

\begin{exo}[Sabit dış basınç altında sıkıştırma]{compression-pression-exterieure-constante}
\seoready{true}
\keywords{basınç kuvvetlerinin işi, sabit dış basınç, sıkıştırma, genleşme, boşluğa genleşme}

Bir gaz, sabit $P_0$ dış basıncı altında $V_A$ hacminden $V_B<V_A$ hacmine sıkıştırılır.

\begin{enumerate}
\item Gazın aldığı işi hesaplayın.
\item Gaz aynı $P_0$ dış basıncına karşı $V_B$'den $V_A$'ya geri döner. Alınan işi hesaplayıp sıkıştırmadaki işle karşılaştırın.
\item Gaz boşluğa genleşirken $V_B$'den $V_A$'ya genleşmeyi yeniden ele alın. Elde edilen üç işaretin anlamını açıklayın.
\end{enumerate}

\begin{indication}
$W=-\int P_{\mathrm{ext}}\,\mathrm{d}V$ bağıntısını kullanın. Süreç yarı-statik
olmadığında bile tümellenmesi gereken basınç \emph{dış} basınçtır.
\end{indication}

\begin{solution}
\textbf{Soru 1.} Dış basınç sabittir; dolayısıyla
\[
W_{A\to B}=-P_0(V_B-V_A)=P_0(V_A-V_B)>0.
\]
Gaz sıkıştırma sırasında iş alır.

\textbf{Soru 2.} Geri dönüş için
\[
W_{B\to A}=-P_0(V_A-V_B)<0.
\]
İki süreç de aynı sabit dış basınca karşı gerçekleştirildiğinden iki iş birbirinin tersidir.

\textbf{Soru 3.} Boşlukta $P_{\mathrm{ext}}=0$ olduğundan
\[
W_{B\to A}^{\mathrm{vide}}=0.
\]
Dolayısıyla genleşme otomatik olarak negatif iş anlamına gelmez: gerçekten
geri itilmiş bir dış kuvvet bulunmalıdır.
\end{solution}
\end{exo}
